% ---------------------------------------------------------------------
% CHAPTER — SYNTHESIS
% Working draft. Prose states the conditioned-answer structure of each RQ.
% "% INSERT" markers show where a specific result / effect size is slotted
% from the named study, so no quantity is asserted here before its chapter
% reports it.
% ---------------------------------------------------------------------
\chapter{Synthesis}
\label{ch:synthesis}
 \chaptertoc
 
This chapter returns to the three research questions of Section~\ref{sec:research-objectives} and states the answer each receives once the studies are read together. The answers are read off the controlled comparisons of Chapter~\ref{ch:methodology}: each study isolates one part of the design space, and the synthesis is what those isolated comparisons say in combination, against the system properties recorded in Table~\ref{table:dataset_details_unified}. None of the three questions resolves to a single ranking of methods. Each resolves to a design choice or a signal together with the conditions under which it carries architectural information, and it is the conditioning, not the ranking, that transfers across systems and across the two tasks.
 
\section{RQ1: Which Design Choices Are Justified}
\label{sec:syn-rq1}
 
RQ1 asks which choices of graph granularity, relation modeling, encoder, and learning objective are justified by measurable benefit. Read across the studies, the answer is that the benefit of a choice is decided by the graph it operates on rather than by the choice in isolation, and the studies make this concrete one axis at a time.
 
The comparison of degree-normalized convolution with attention illustrates the point. In supervised mapping, attention is the stronger and safer default: the relation-typed attention encoder is the best-performing learned model, whereas the convolutional encoder does not consistently improve on the graph-free baseline, and the depth at which it is useful depends on the granularity, connectivity, and homophily of the system rather than being fixed. % INSERT: C2A ICSA main results (GAT best; GCN not consistently > MLP) + depth ablation effect sizes; state the systems where GCN depth helps vs hurts.
The recovery studies compare the same two encoders within the reconstruction framework, where the decision is made without labels. % INSERT: GAER GCN vs GAT verdict + the condition under which one is preferred; note granularity (number of clusters) effect.
 
Keeping dependency types apart through relation-specific parameters, rather than collapsing them into a single weighted graph, is justified where the types carry distinct architectural meaning and does not pay for itself where they do not. % INSERT: C2A typed-vs-collapsed and edge-direction result; GAER heterogeneous-encoder motivation and any type-specific gain.
The learning objective is the axis particular to unsupervised recovery, where a reconstruction objective and a contrastive objective are available in the absence of labels. % INSERT: ARL (contrastive) vs GAER (reconstruction) comparison and the condition that separates them.
 
The richest representation extends the graph itself, from a file-only graph to one that adds folder and member-level nodes, processed by a meta-relation attention encoder. Whether this added structure is justified is not settled at the level of the encoder but at the level of what the extra nodes contribute, which is the subject of RQ3. % INSERT: HetMap HGT vs GAT/MLP headline (RQ1 table) — state where the heterogeneous representation wins and where it does not.
 
\section{RQ2: What Each Input Signal Carries}
\label{sec:syn-rq2}
 
RQ2 asks what architectural information dependency structure, project organization, and code identifiers carry, and under which system properties each is reliable. The answer is stated per signal, and each signal's reliability is tied to a property that Table~\ref{table:dataset_details_unified} records or that the studies measure, so that the claim is conditional rather than general.
 
The contribution of the dependency graph is isolated by the graph-free baseline, which predicts module assignments from node features alone under the same training procedure. What the graph adds over the features is therefore measured directly rather than inferred. % INSERT: C2A ICSA RQ2 result — graph contribution beyond features, and its dependence on homophily / dependency cohesion (name the systems where the graph helps and where it does not). Coverage effect (MLP maps ~60% at 0.95 confidence vs ~97% with the graph) belongs here.
Read against homophily, the dependency signal is informative where dependencies concentrate within reference modules and weak where cross-module coupling, for example through widely used utility code, dominates the local neighborhood.
 
Project organization enters as folder structure, and its reliability tracks how closely the folder hierarchy aligns with the reference decomposition. Where that alignment is high, folder structure is a strong architectural proxy and can carry much of the result on its own; where it is low, it is noise, since utility directories mix entities from several modules. % INSERT: GAER dependency-only vs folder-augmented comparison; StructSAR folder-organization result. The "shown by removing it" ablation for directory structure (contribution C4) is stated here.
Code identifiers carry domain vocabulary that aligns with module membership where the vocabulary is separable across modules, and they blur where the same terms recur across modules. % INSERT: ARL lexical-signal result; HetMap file-feature (Word2Vec / node2vec) behavior as it bears on lexical vs structural initialization.
 
\section{RQ3: How Signals Interact in a Richer Graph}
\label{sec:syn-rq3}
 
RQ3 asks how the usefulness of structural, lexical, and contextual file representations changes once they are available together within a heterogeneous graph, and where the sources are complementary rather than redundant. It is answered principally by the capstone study, which varies file representation and graph structure independently while holding the learning setting fixed, and it functions as a test of whether the RQ1 and RQ2 conclusions survive once the signals are combined.
 
The distinction the study draws is between a gain that reflects new architectural evidence and a gain that reflects a signal already aligned with the reference. Adding member-level structure to the folder-file graph improves mapping where the added interactions separate files that the coarser graph confused, and does not improve it where folder organization has already saturated the achievable performance. % INSERT: HetMap RQ3 significance table (F vs MLP, FM vs MLP, FM vs F) read by system; name the systems where member structure adds (complementary) and where folder alignment saturates (redundant).
The interaction also depends on how the file node is initialized. A lexical initialization and a structural, random-walk initialization respond differently to the same added graph structure, because a representation that already encodes neighborhood structure has less to gain, and can lose, from further message passing over that structure. % INSERT: HetMap Word2Vec vs node2vec contrast from the RQ3 table — state where message passing helps the lexical init and where it is neutral or harmful for the structural init.
 
The consequence for the thesis is that the value of a signal is not additive. A signal that is informative in isolation, as established under RQ2, can be redundant once another aligned signal is present, and the heterogeneous model makes this visible rather than hiding it inside an aggregate score.
 
\section{Consolidated Contributions}
\label{sec:syn-contributions}
 
The three conditioned answers restate the contributions claimed in Chapter~\ref{ch:introduction}. That learned representations outperform hand-tuned fusion, most clearly at the scale where fixed weightings are most brittle, is the combined verdict of RQ1 and RQ2 across the two tasks. The isolation of what each signal contributes and where it fails is the substance of RQ2. That the same design choices and signals behave consistently across unsupervised recovery and supervised mapping, which no single study can claim, is what the shared design space of Chapter~\ref{ch:methodology} makes demonstrable. And that directory structure is a strong architectural proxy, shown by removing it rather than by asserting it, is the specific result folded into RQ2.